\documentclass[aspectratio=169, 10pt]{beamer}

% ── Theme & colors ────────────────────────────────────────────────────────────
\usetheme{Madrid}
\usecolortheme{default}
\definecolor{primary}{RGB}{0,84,166}
\definecolor{accent}{RGB}{255,140,0}
\definecolor{lightgray}{RGB}{245,245,245}
\definecolor{darkgreen}{RGB}{0,128,64}
\definecolor{darkred}{RGB}{180,0,0}
\definecolor{whitesoft}{RGB}{220,220,220}

\setbeamercolor{palette primary}{bg=primary,fg=white}
\setbeamercolor{palette secondary}{bg=primary!80,fg=white}
\setbeamercolor{palette tertiary}{bg=primary!60,fg=white}
\setbeamercolor{palette quaternary}{bg=primary,fg=white}
\setbeamercolor{frametitle}{bg=primary,fg=white}
\setbeamercolor{title}{fg=white}
\setbeamercolor{block title}{bg=primary,fg=white}
\setbeamercolor{block title alerted}{bg=accent,fg=white}
\setbeamercolor{alerted text}{fg=accent}
\setbeamercolor{structure}{fg=primary}

% ── Packages ──────────────────────────────────────────────────────────────────
\usepackage[utf8]{inputenc}
\usepackage[T1]{fontenc}
\usepackage[french]{babel}
\usepackage{booktabs}
\usepackage{array}
\usepackage{xcolor}
\usepackage{tikz}
\usetikzlibrary{
  shapes.geometric,
  arrows.meta,
  positioning,
  fit,
  backgrounds,
  decorations.pathreplacing
}
\usepackage{listings}
\usepackage{multirow}
\usepackage{colortbl}
\usepackage{graphicx}
\usepackage{textcomp}

% ── Listings style ────────────────────────────────────────────────────────────
\lstset{
  basicstyle=\ttfamily\tiny,
  backgroundcolor=\color{lightgray},
  frame=single,
  framerule=0.4pt,
  breaklines=true,
  columns=fullflexible
}

% ── TikZ styles ───────────────────────────────────────────────────────────────
\tikzset{
  arrow/.style={-{Stealth[length=6pt]}, thick, color=primary},
  box/.style={
    rectangle, rounded corners=4pt, draw=primary, thick,
    fill=primary!10, text width=2.8cm, align=center,
    minimum height=0.85cm, font=\small\bfseries
  },
  boxsmall/.style={
    rectangle, rounded corners=3pt, draw=primary, thick,
    fill=primary!8, text width=2.2cm, align=center,
    minimum height=0.65cm, font=\scriptsize
  },
  oldbox/.style={
    rectangle, rounded corners=3pt, draw=darkred, thick,
    fill=darkred!8, text width=3.5cm, align=center,
    minimum height=0.65cm, font=\scriptsize
  },
  newbox/.style={
    rectangle, rounded corners=3pt, draw=darkgreen, thick,
    fill=darkgreen!8, text width=3.5cm, align=center,
    minimum height=0.65cm, font=\scriptsize
  },
}

% ── Metadata ──────────────────────────────────────────────────────────────────
\title[RAG Hybride -- Extraction Financiere]{%
  \textbf{Systeme RAG Hybride}\\[4pt]
  \large Extraction Automatique de Donnees Financieres\\
  depuis les Rapports AMMC
}
\author[Presentation]{%
  \textbf{Nom Prenom}\\[4pt]
  \small Encadrant~: Pr.~Nom Encadrant
}
\institute[]{Filiere -- Etablissement}
\date{Avril 2026}

% ══════════════════════════════════════════════════════════════════════════════
\begin{document}
% ══════════════════════════════════════════════════════════════════════════════

% ── Title ─────────────────────────────────────────────────────────────────────
\begin{frame}[plain]
  \begin{tikzpicture}[overlay, remember picture]
    \fill[primary] (current page.north west) rectangle
                   ([yshift=-3.6cm]current page.north east);
    \fill[accent]  ([yshift=-3.6cm]current page.north west) rectangle
                   ([yshift=-3.85cm]current page.north east);
  \end{tikzpicture}
  \vspace{0.9cm}
  \titlepage
\end{frame}

% ── Plan ──────────────────────────────────────────────────────────────────────
\begin{frame}{Plan de la presentation}
  \tableofcontents[hideallsubsections]
\end{frame}

% ══════════════════════════════════════════════════════════════════════════════
\section{Contexte et Problematique}
% ══════════════════════════════════════════════════════════════════════════════

\begin{frame}{Contexte et Problematique}
  \begin{columns}[T]
    \column{0.52\textwidth}
      \begin{alertblock}{Probleme constate}
        \begin{itemize}
          \item L'AMMC publie des centaines de \textbf{rapports financiers PDF}
                chaque annee (bilans, CPC, flux de tresorerie\ldots)
          \item L'extraction manuelle est \textbf{longue, couteuse et sujette
                aux erreurs}
          \item Les PDFs sont heterogenes :
                \begin{itemize}\small
                  \item Documents natifs vs \textbf{scannes}
                  \item Pages multi-tableaux
                  \item Normes \textbf{CGNC} et \textbf{PCEC} (secteur bancaire)
                \end{itemize}
        \end{itemize}
      \end{alertblock}

    \column{0.45\textwidth}
      \begin{block}{Objectif du projet}
        Construire un systeme \textbf{RAG Hybride} capable de :
        \begin{enumerate}\small
          \item Ingerer automatiquement les PDFs
          \item Localiser les tableaux par \textbf{recherche semantique}
          \item Extraire les donnees via une \textbf{cascade LLM + methodes classiques}
          \item Valider et exporter en \textbf{Excel structure}
        \end{enumerate}
      \end{block}
      \vspace{0.3cm}
      \begin{tikzpicture}
        \node[
          fill=accent!20, draw=accent, rounded corners=4pt,
          font=\small, text width=5.4cm, align=center, inner sep=7pt
        ]{
          Emetteurs cibles : Attijariwafa, BMCI,\\Addoha, AGMA\ldots
        };
      \end{tikzpicture}
  \end{columns}
\end{frame}

% ══════════════════════════════════════════════════════════════════════════════
\section{Architecture Generale}
% ══════════════════════════════════════════════════════════════════════════════

\begin{frame}{Architecture Generale du Systeme}
  \centering
  \begin{tikzpicture}[node distance=0.5cm and 0.95cm, scale=0.87, transform shape]
    \node[box, fill=accent!20, draw=accent] (sc)  {Scraper AMMC\\{\tiny Download PDF}};
    \node[box, right=of sc]                (ing) {Ingestion\\{\tiny Chunks + Embed.}};
    \node[box, right=of ing]               (idx) {Index FAISS\\{\tiny Vecteurs}};
    \node[box, right=of idx]               (ret) {Retrieval\\{\tiny Rech. semantique}};
    \node[box, right=of ret, fill=darkgreen!15, draw=darkgreen]
                                           (ext) {Extraction\\{\tiny Cascade LLM}};
    \node[box, right=of ext, fill=accent!20, draw=accent]
                                           (exp) {Export Excel\\{\tiny .xlsx formate}};

    \draw[arrow] (sc)--(ing);
    \draw[arrow] (ing)--(idx);
    \draw[arrow] (idx)--(ret);
    \draw[arrow] (ret)--(ext);
    \draw[arrow] (ext)--(exp);

    \node[font=\tiny, color=gray, below=0.22cm of sc]  {Phase 0};
    \node[font=\tiny, color=gray, below=0.22cm of ing] {Phase 1};
    \node[font=\tiny, color=gray, below=0.22cm of idx] {Phase 1};
    \node[font=\tiny, color=gray, below=0.22cm of ret] {Phase 2};
    \node[font=\tiny, color=gray, below=0.22cm of ext] {Phase 3};
    \node[font=\tiny, color=gray, below=0.22cm of exp] {Phase 4};

    \draw[
      decorate,
      decoration={brace, mirror, amplitude=6pt},
      thick, color=primary
    ]
    ([yshift=-0.75cm]ing.south west) -- ([yshift=-0.75cm]ret.south east)
    node[midway, below=9pt, font=\small\bfseries, color=primary]{RAG Pipeline};
  \end{tikzpicture}

  \vspace{0.55cm}
  \begin{columns}[T]
    \column{0.23\textwidth}
      \centering\small\textbf{[DB] Donnees}\\
      \tiny pdfplumber $\cdot$ PyMuPDF
    \column{0.23\textwidth}
      \centering\small\textbf{[~] Embedding}\\
      \tiny all-MiniLM-L6-v2
    \column{0.23\textwidth}
      \centering\small\textbf{[Q] Retrieval}\\
      \tiny FAISS cosine similarity
    \column{0.23\textwidth}
      \centering\small\textbf{[AI] LLM}\\
      \tiny Groq $\cdot$ Gemini $\cdot$ OpenAI
  \end{columns}
\end{frame}

% ══════════════════════════════════════════════════════════════════════════════
\section{Etapes du Pipeline}
% ══════════════════════════════════════════════════════════════════════════════

\begin{frame}{Etape 1 --- Scraping et Ingestion}
  \begin{columns}[T]
    \column{0.48\textwidth}
      \begin{block}{Scraper AMMC (\texttt{scraper.py})}
        \textbf{Strategie de recuperation (priorite)} :
        \begin{enumerate}\small
          \item Cache local \texttt{.pdf}
          \item Overrides manuels (\texttt{pdf\_overrides.json})
          \item Parsing du formulaire AMMC (Emetteur + Annee)
          \item Fallback page de la liste des emetteurs
        \end{enumerate}
      \end{block}

    \column{0.48\textwidth}
      \begin{block}{Ingestion (\texttt{ingest.py})}
        \textbf{Deux modes d'extraction} :
        \begin{itemize}\small
          \item \textbf{Full} (pdfplumber) --- preserve les tableaux Markdown
          \item \textbf{Fast} (PyMuPDF) --- 5--10$\times$ plus rapide
        \end{itemize}
        \textbf{Sortie} :
        \begin{itemize}\small
          \item Chunks de 512 caract. + tableaux complets
          \item Metadata : \{page, kind, section\_hint\}
          \item \texttt{embeddings.npy} + \texttt{metadata.json}
        \end{itemize}
      \end{block}
  \end{columns}

  \vspace{0.3cm}
  \begin{tikzpicture}[scale=0.85, transform shape, node distance=0.45cm and 0.75cm]
    \node[boxsmall, fill=accent!15, draw=accent] (pdf)   {PDF AMMC};
    \node[boxsmall, right=of pdf]                (chk)   {Chunking\\{\tiny 512 chars}};
    \node[boxsmall, right=of chk]                (emb)   {Embedding\\{\tiny MiniLM-L6}};
    \node[boxsmall, right=of emb]                (fss)   {Index FAISS\\{\tiny .npy + .json}};
    \node[boxsmall, right=of fss, fill=darkgreen!15, draw=darkgreen]
                                                 (ok)    {Pret pour\\Retrieval};
    \draw[arrow] (pdf)--(chk);
    \draw[arrow] (chk)--(emb);
    \draw[arrow] (emb)--(fss);
    \draw[arrow] (fss)--(ok);
  \end{tikzpicture}

  \vspace{0.2cm}
  \small\textbf{Detection de section} : 23 regles par frequence de mots-cles
  (BILAN ACTIF, BILAN PASSIF, CPC --- normes CGNC \& PCEC)
\end{frame}

%──────────────────────────────────────────────────────────────────────────────

\begin{frame}{Etape 2 --- Recherche Semantique (Retrieval)}
  \begin{columns}[T]
    \column{0.52\textwidth}
      \begin{block}{SimpleVectorRetriever (\texttt{retriever.py})}
        \begin{enumerate}\small
          \item Encodage de la requete avec \textbf{MiniLM-L6-v2}
          \item \textbf{Similarite cosinus} sur l'index FAISS
          \item Retourne les \textbf{top-K chunks} (defaut K=8)
                avec score de pertinence
          \item Filtrage optionnel par \texttt{kind} : text | table
        \end{enumerate}
      \end{block}

      \begin{alertblock}{Classification d'intention}
        Heuristique sur les mots-cles de la question :\\
        \small $\rightarrow$ (intent, table\_type)\\
        \textit{ex. : ``table\_reconstruction'', ``bilan\_actif''}
      \end{alertblock}

    \column{0.44\textwidth}
      \centering
      \begin{tikzpicture}[scale=0.80, transform shape, node distance=0.45cm]
        \node[box, fill=accent!15, draw=accent, text width=3.5cm]
              (q)   {``Bilan actif 2024 ?''};
        \node[box, below=of q, text width=3.5cm]
              (enc) {Encodage\\{\tiny MiniLM-L6-v2}};
        \node[box, below=of enc, text width=3.5cm]
              (sim) {Similarite cosinus\\{\tiny FAISS}};
        \node[box, below=of sim, fill=darkgreen!15, draw=darkgreen, text width=3.5cm]
              (res) {Top-8 chunks\\{\tiny + section\_hint}};
        \draw[arrow] (q)  --(enc);
        \draw[arrow] (enc)--(sim);
        \draw[arrow] (sim)--(res);
      \end{tikzpicture}
  \end{columns}
\end{frame}

% ══════════════════════════════════════════════════════════════════════════════
\section{Module d'Extraction}
% ══════════════════════════════════════════════════════════════════════════════

\begin{frame}{Etape 3 --- Extraction : Ancienne vs Nouvelle Methode}
  \begin{columns}[T]

    % ── Ancienne ─────────────────────────────────────────────────────────────
    \column{0.46\textwidth}
      \begin{center}
        \colorbox{darkred!15}{%
          \parbox{5.0cm}{%
            \centering\small\bfseries\color{darkred}
            [Avant] Ancienne Methode
          }%
        }
      \end{center}
      \vspace{0.1cm}
      \begin{tikzpicture}[node distance=0.28cm, scale=0.82, transform shape]
        \node[oldbox] (v1)
          {Niveau 1 --- Vision LLM\\
           \tiny Rendu page $\rightarrow$ image $\rightarrow$ LLM};
        \node[oldbox, below=of v1] (v2)
          {Niveau 2 --- JSON structure\\
           \tiny LLM $\rightarrow$ JSON (parsing manuel)};
        \node[oldbox, below=of v2] (v3)
          {Niveau 3 --- Markdown\\
           \tiny LLM $\rightarrow$ Markdown $\rightarrow$ RegEx};
        \node[oldbox, below=of v3] (v4)
          {Niveau 4 --- Fallback\\
           \tiny Cross-provider (multi-LLM)};
        \node[oldbox, below=of v4, fill=orange!15, draw=orange] (mt)
          {+ Detection multi-tableaux\\
           \tiny Segmentation par zones};
        \node[oldbox, below=of mt, fill=orange!15, draw=orange] (cc)
          {+ Controle de completude\\
           \tiny Score 0.0 -- 1.0};
        \draw[arrow,color=darkred] (v1)--(v2);
        \draw[arrow,color=darkred] (v2)--(v3);
        \draw[arrow,color=darkred] (v3)--(v4);
        \draw[arrow,color=darkred] (v4)--(mt);
        \draw[arrow,color=darkred] (mt)--(cc);
      \end{tikzpicture}

    % ── Separator ────────────────────────────────────────────────────────────
    \column{0.05\textwidth}
      \vspace{2.2cm}
      \centering\textcolor{gray}{\Large $\Rightarrow$}

    % ── Nouvelle ─────────────────────────────────────────────────────────────
    \column{0.46\textwidth}
      \begin{center}
        \colorbox{darkgreen!15}{%
          \parbox{5.0cm}{%
            \centering\small\bfseries\color{darkgreen}
            [2025] Nouvelle Methode
          }%
        }
      \end{center}
      \vspace{0.1cm}
      \begin{tikzpicture}[node distance=0.28cm, scale=0.82, transform shape]
        \node[newbox] (n1)
          {Tier 1 --- Vision LLM\\
           \tiny Groq Llama-4 / Gemini 2.0 Flash\\
           \tiny PDF $\rightarrow$ 350 DPI $\rightarrow$ Base64 $\rightarrow$ JSON};
        \node[newbox, below=of n1] (n2)
          {Tier 2 --- Structured Output\\
           \tiny pymupdf4llm $\rightarrow$ Markdown\\
           \tiny Groq 70B + \textbf{response\_format=json}};
        \node[newbox, below=of n2] (n3)
          {Tier 3 --- Markdown Extract\\
           \tiny LLM Markdown $\rightarrow$ RegEx $\rightarrow$ DataFrame};
        \node[newbox, below=of n3] (n4)
          {Tier 4 --- Methodes classiques\\
           \tiny Camelot (lattice/stream)\\
           \tiny pdfplumber $\rightarrow$ clustering X/Y};
        \draw[arrow,color=darkgreen] (n1)--(n2);
        \draw[arrow,color=darkgreen] (n2)--(n3);
        \draw[arrow,color=darkgreen] (n3)--(n4);
      \end{tikzpicture}
  \end{columns}
\end{frame}

%──────────────────────────────────────────────────────────────────────────────

\begin{frame}{Ameliorations Cles de la Nouvelle Methode}
  \begin{columns}[T]
    \column{0.48\textwidth}
      \begin{block}{[+] Structured JSON Output}
        \begin{itemize}\small
          \item \textbf{Avant} : Parsing manuel du JSON LLM
                $\Rightarrow$ erreurs frequentes
          \item \textbf{Apres} : \texttt{response\_format=json\_object}
                (Groq) $\Rightarrow$ \textbf{zero erreur de parsing}
        \end{itemize}
      \end{block}

      \begin{block}{[+] pymupdf4llm}
        \begin{itemize}\small
          \item \textbf{Avant} : pdfplumber brut
                $\Rightarrow$ perte de structure
          \item \textbf{Apres} : Markdown fidele a la mise en page
                $\Rightarrow$ meilleure comprehension LLM
        \end{itemize}
      \end{block}

    \column{0.48\textwidth}
      \begin{block}{[+] Modeles Vision}
        \begin{itemize}\small
          \item Llama-4 Maverick (Groq) + Gemini 2.0 Flash
          \item Gerent les PDFs \textbf{scannes et natifs}
          \item Selection automatique selon qualite du contenu
        \end{itemize}
      \end{block}

      \begin{block}{[+] Fusion de colonnes}
        \begin{itemize}\small
          \item Corrige les nombres fragmentes par le separateur
                de milliers
          \item Ex. : ``13 157'' + ``683'' $\rightarrow$ ``13 157 683''
          \item Heuristique : $>50\%$ de fragments $\Rightarrow$ fusion
        \end{itemize}
      \end{block}
  \end{columns}
\end{frame}

%──────────────────────────────────────────────────────────────────────────────

\begin{frame}{Modules de Validation et Controle Qualite}
  \begin{columns}[T]
    \column{0.48\textwidth}
      \begin{block}{Completude (\texttt{completeness\_checker.py})}
        \begin{itemize}\small
          \item Verifie la presence des \textbf{ancres critiques}
                (ex. ``TOTAL DE L'ACTIF'')
          \item Valide le nombre minimal de lignes/colonnes
          \item Genere un \textbf{score 0.0 -- 1.0}
          \item Decide si un tier superieur est necessaire
        \end{itemize}
      \end{block}

      \begin{block}{Multi-Tableaux (\texttt{multi\_table\_detector.py})}
        \begin{itemize}\small
          \item Detecte 2--3 tableaux sur une meme page
          \item Extraction par \textbf{zones bornees}
                (mots-cles + coordonnees)
          \item Evite l'extraction partielle
        \end{itemize}
      \end{block}

    \column{0.48\textwidth}
      \begin{block}{Extraction Positionnelle (\texttt{positional.py})}
        \begin{itemize}\small
          \item Pour tableaux sans bordures visibles
          \item Clustering \textbf{X/Y} des cellules
          \item 3 strategies en cascade :
                \begin{enumerate}\scriptsize
                  \item pdfplumber natif
                  \item Clustering X/Y
                  \item Separation par espaces
                \end{enumerate}
        \end{itemize}
      \end{block}

      \begin{alertblock}{Resultat ExtractionAgent}
        \texttt{ExtractionResult} contient :\\[3pt]
        \texttt{df}, \texttt{method\_used},\\
        \texttt{confidence}, \texttt{completeness\_score}
      \end{alertblock}
  \end{columns}
\end{frame}

% ══════════════════════════════════════════════════════════════════════════════
\section{Question-Reponse Semantique}
% ══════════════════════════════════════════════════════════════════════════════

\begin{frame}{Pipeline Question-Reponse Semantique}
  \centering
  \begin{tikzpicture}[node distance=0.38cm and 0.9cm, scale=0.84, transform shape]
    \node[box, fill=accent!15, draw=accent] (q)   {Question\\utilisateur};
    \node[box, right=of q]                  (cls) {Classification\\d'intention};
    \node[box, right=of cls]                (ret) {Retrieval\\top-8 chunks};
    \node[box, right=of ret]                (ctx) {Formatage\\contexte};
    \node[box, right=of ctx, fill=darkgreen!15, draw=darkgreen]
                                            (llm) {LLM\\GPT / Groq};
    \node[box, right=of llm, fill=accent!15, draw=accent]
                                            (ans) {Reponse\\Markdown};
    \draw[arrow] (q)--(cls);
    \draw[arrow] (cls)--(ret);
    \draw[arrow] (ret)--(ctx);
    \draw[arrow] (ctx)--(llm);
    \draw[arrow] (llm)--(ans);
  \end{tikzpicture}

  \vspace{0.45cm}
  \begin{columns}[T]
    \column{0.48\textwidth}
      \begin{block}{Prompting intelligent}
        \begin{itemize}\small
          \item Intent \textbf{table\_reconstruction}
                $\Rightarrow$ instruction LLM pour construire
                un tableau Markdown
          \item Adapte par type : BILAN (Actif/Passif),
                CPC (Produits/Charges), Flux\ldots
          \item Consigne : \textit{``Ne pas inventer de chiffres''}
        \end{itemize}
      \end{block}

    \column{0.48\textwidth}
      \begin{block}{Section hints dans le contexte}
        \begin{itemize}\small
          \item Chaque chunk embarque un \textbf{section\_hint}
                (type de tableau detecte a l'ingestion)
          \item Aide le LLM a distinguer Actif / Passif / CPC
          \item Reduit les confusions entre tableaux similaires
        \end{itemize}
      \end{block}
  \end{columns}
\end{frame}

% ══════════════════════════════════════════════════════════════════════════════
\section{API et Deploiement}
% ══════════════════════════════════════════════════════════════════════════════

\begin{frame}{API REST et Interface (\texttt{api.py})}
  \begin{columns}[T]
    \column{0.48\textwidth}
      \begin{block}{FastAPI Service}
        \begin{itemize}\small
          \item CORS active (frontend Vite/React)
          \item \textbf{File d'attente asynchrone} pour l'extraction
                (ThreadPoolExecutor)
          \item Endpoints principaux :
                \begin{itemize}\scriptsize
                  \item \texttt{POST /ask} --- Q\&R semantique
                  \item \texttt{POST /extract} --- Extraction tableau
                  \item \texttt{GET /export\_xlsx} --- Export Excel
                  \item \texttt{POST /ingest} --- Indexation PDF
                \end{itemize}
        \end{itemize}
      \end{block}

    \column{0.48\textwidth}
      \begin{block}{Configuration multi-provider}
        \begin{itemize}\small
          \item \textbf{Groq} (Llama-4, Llama-3.3-70B) --- rapide
          \item \textbf{Gemini 2.0 Flash} --- vision scans
          \item \textbf{OpenAI GPT} --- Q\&R narrative
          \item Selection automatique selon le tier actif
        \end{itemize}
      \end{block}

      \begin{block}{Demarrage}
        \begin{lstlisting}
uvicorn api:app --reload --port 8000
        \end{lstlisting}
        \small Supporte : AGMA, Attijariwafa,\\Addoha, BMCI\ldots
      \end{block}
  \end{columns}
\end{frame}

% ══════════════════════════════════════════════════════════════════════════════
\section{Technologies et Resultats}
% ══════════════════════════════════════════════════════════════════════════════

\begin{frame}{Stack Technologique}
  \centering
  \vspace{0.15cm}
  \begin{tabular}{@{}lll@{}}
    \toprule
    \textbf{Couche}      & \textbf{Technologie}                   & \textbf{Role} \\
    \midrule
    Extraction PDF       & pdfplumber, PyMuPDF, pymupdf4llm       & Texte + tableaux \\
    Vision OCR           & Groq Llama-4, Gemini 2.0 Flash         & PDFs scannes \\
    LLM Structure        & Groq llama-3.3-70B (JSON mode)         & Extraction fiable \\
    Embeddings           & sentence-transformers/all-MiniLM-L6-v2 & Vecteurs semantiques \\
    Base vectorielle     & FAISS (CPU)                            & Recherche rapide \\
    Fallback classique   & Camelot, Tabula, pdfplumber            & Grilles detectees \\
    API                  & FastAPI + uvicorn                      & Service REST \\
    Export               & Pandas + xlsxwriter                    & Excel formate \\
    \bottomrule
  \end{tabular}

  \vspace{0.45cm}
  \begin{columns}[T]
    \column{0.3\textwidth}
      \begin{tikzpicture}
        \node[fill=primary!10, draw=primary, rounded corners=4pt,
              text width=3.5cm, align=center, inner sep=7pt, font=\small]{
          \textbf{Ingestion}\\[4pt]
          Fast mode :\\1--2 sec/page\\
          Full mode :\\5--10 sec/page
        };
      \end{tikzpicture}
    \column{0.3\textwidth}
      \begin{tikzpicture}
        \node[fill=darkgreen!10, draw=darkgreen, rounded corners=4pt,
              text width=3.5cm, align=center, inner sep=7pt, font=\small]{
          \textbf{Extraction}\\[4pt]
          Vision LLM :\\2--3 sec/page\\
          Cascade auto\\si necessaire
        };
      \end{tikzpicture}
    \column{0.3\textwidth}
      \begin{tikzpicture}
        \node[fill=accent!15, draw=accent, rounded corners=4pt,
              text width=3.5cm, align=center, inner sep=7pt, font=\small]{
          \textbf{Scalabilite}\\[4pt]
          Workers paralleles\\multi-PDF\\Jobs asynchrones\\via API
        };
      \end{tikzpicture}
  \end{columns}
\end{frame}

% ══════════════════════════════════════════════════════════════════════════════
\section{Conclusion et Perspectives}
% ══════════════════════════════════════════════════════════════════════════════

\begin{frame}{Conclusion et Perspectives}
  \begin{columns}[T]
    \column{0.50\textwidth}
      \begin{block}{Bilan du projet}
        \begin{itemize}\small
          \item Systeme \textbf{end-to-end} operationnel :\\
                scraping $\to$ ingestion $\to$ retrieval
                $\to$ extraction $\to$ export
          \item Gestion des cas difficiles : scans,
                multi-tableaux, normes CGNC/PCEC
          \item \textbf{Cascade robuste} : 4 tiers d'extraction
                avec validation automatique
          \item Score de completude par tableau extrait
        \end{itemize}
      \end{block}

    \column{0.46\textwidth}
      \begin{alertblock}{Perspectives}
        \begin{itemize}\small
          \item Extension a d'autres emetteurs AMMC
          \item Fine-tuning sur documents financiers marocains
          \item Interface front-end React complete
          \item Base de donnees pour historique multi-annees
          \item Detection automatique d'anomalies comptables
        \end{itemize}
      \end{alertblock}
  \end{columns}

  \vspace{0.4cm}
  \centering
  \begin{tikzpicture}
    \node[
      fill=primary!10, draw=primary, rounded corners=5pt,
      text width=11.5cm, align=center, inner sep=10pt
    ]{
      \textbf{Apport principal :} remplacer des heures d'extraction manuelle\\
      par un pipeline automatise, fiable et extensible, combinant\\
      \textbf{recherche semantique (RAG)} et \textbf{LLMs multimodaux}.
    };
  \end{tikzpicture}
\end{frame}

% ── Slide finale ──────────────────────────────────────────────────────────────
\begin{frame}[plain]
  \begin{tikzpicture}[overlay, remember picture]
    \fill[primary] (current page.north west) rectangle (current page.south east);
  \end{tikzpicture}
  \centering
  \vspace{2.5cm}
  {\color{white}\Huge\bfseries Merci pour votre attention}\\[1.2cm]
  {\color{whitesoft}\Large Questions ?}\\[1.5cm]
  {\color{accent}\large votre.email@domaine.ma}
\end{frame}

\end{document}
